\documentclass[11pt,a4paper]{article}
\makeatletter\def\input@path{{../../templates/}}\makeatother
\usepackage{avaliacao-poo}
\configurarAvaliacao{Checkpoint 1}{Programação Orientada a Objetos}
\validarCheckpointDuasPaginas
\begin{document}
\cabecalhoAvaliacao
\blocoQuadroRespostas
\cabecalhoCenario{Elevador}
Os arquivos \codigoInline{Elevador.java} e \codigoInline{Main.java} estão no mesmo pacote.
Considere exatamente o código abaixo.

\noindent\begin{minipage}[t]{.37\textwidth}
\vspace{0pt}
\begin{painelreferencia}{Elevador.java}
\begin{lstlisting}[style=cenariocodigo]
public class Elevador {
    int andarAtual = 0;
    boolean portaAberta = false;

    void subir() {
        andarAtual = andarAtual + 1;
    }
    void abrirPorta() {
        portaAberta = true;
    }
    void fecharPorta() {
        portaAberta = false;
    }
    int getAndarAtual() {
        return andarAtual;
    }
    boolean getPortaAberta() {
        return portaAberta;
    }
}
\end{lstlisting}
\end{painelreferencia}
\end{minipage}\hfill
\begin{minipage}[t]{.61\textwidth}
\vspace{0pt}
\begin{painelreferencia}{Main.java}
\begin{lstlisting}[style=cenariocodigo]
public class Main {
    public static void main(String[] args) {
        Elevador principal = new Elevador();
        Elevador painel = principal;
        Elevador servico = new Elevador();

        principal.subir();
        painel.abrirPorta();

        System.out.println(principal.getAndarAtual());
        System.out.println(principal.getPortaAberta());

        painel.fecharPorta();
        painel.subir();

        System.out.println(painel.getAndarAtual());
        System.out.println(painel.getPortaAberta());

        servico.subir();
        servico.abrirPorta();

        System.out.println(servico.getAndarAtual());
    }
}
\end{lstlisting}
\end{painelreferencia}
\end{minipage}
\inicioQuestoes

\questaoDireta{1}{Faça o teste de mesa e indique o que será impresso em cada linha.}{25 pontos}
\itensDuasColunas{
  \itemSimples{a}{Linha 10}
  \itemSimples{b}{Linha 11}
  \itemSimples{c}{Linha 16}
}{
  \itemSimples{d}{Linha 17}
  \itemSimples{e}{Linha 22}
}

\questaoDireta{2}{Considere a evolução da classe e marque cada afirmação como V ou F.}{25 pontos}
Os getters e os métodos \codigoInline{abrirPorta()} e \codigoInline{fecharPorta()} permanecem inalterados.
\begin{painelEvolucao}{Trechos alterados em Elevador.java}
\begin{lstlisting}[style=evolucaocodigo]
private int andarAtual;
private boolean portaAberta;
void subir() {
    if (portaAberta == false) {
        if (andarAtual < 3) {
            andarAtual = andarAtual + 1;
        }
    }
}
\end{lstlisting}
\end{painelEvolucao}

\itemVF{a}{Em \codigoInline{Main}, a instrução \codigoInline{principal.andarAtual = 2;} compila.}
\itemVF{b}{O método \codigoInline{subir()} pode alterar \codigoInline{andarAtual}, embora o campo seja privado.}
\itemVF{c}{Do térreo, após \codigoInline{abrirPorta()} e \codigoInline{subir()}, o andar continua em \codigoInline{0}.}
\itemVF{d}{Com a porta fechada e o andar em \codigoInline{3}, chamar \codigoInline{subir()} leva ao andar \codigoInline{4}.}
\itemVF{e}{Do térreo, após \codigoInline{abrirPorta()}, \codigoInline{fecharPorta()} e \codigoInline{subir()}, o andar continua em \codigoInline{0}.}

\noindent\begin{minipage}[t]{.485\textwidth}
\questaoDireta{3}{O programa precisa guardar o andar e a situação da porta e executar as ações de subir, abrir e fechar. Qual organização representa melhor um elevador?}{5 pontos}
\begin{alternativas}
\item Manter os dois dados em \codigoInline{Main} e repetir ali todas as ações.
\item Criar uma classe apenas com o andar e a porta e programar todas as ações em \codigoInline{Main}.
\item Criar \codigoInline{Elevador} com o andar, a porta e os métodos que realizam essas ações.
\item Manter os dados em variáveis de \codigoInline{Main} e passar todos eles a cada método externo.
\item Criar uma classe para o andar, outra para a porta e realizar as ações em \codigoInline{Main}.
\end{alternativas}
\questaoDireta{4}{Ao executar \codigoInline{Elevador principal = new Elevador();}, o que cada parte da instrução representa?}{5 pontos}
\begin{alternativas}
\item \codigoInline{Elevador} é a variável; \codigoInline{principal} é a classe; \codigoInline{new} copia um objeto.
\item \codigoInline{Elevador} é a classe; \codigoInline{principal} guarda uma referência para o novo objeto.
\item \codigoInline{principal} é o objeto; \codigoInline{Elevador} guarda a referência usada por \codigoInline{new}.
\item \codigoInline{new Elevador()} é a variável; \codigoInline{principal} declara outra classe.
\item Os três nomes representam o mesmo papel e identificam diretamente o objeto.
\end{alternativas}
\questaoDireta{5}{Quando \codigoInline{abrirPorta()} é executado, o valor de \codigoInline{portaAberta} muda. Qual alternativa descreve corretamente esses dois membros?}{5 pontos}
\begin{alternativas}
\item \codigoInline{portaAberta} é estado; \codigoInline{abrirPorta()} é comportamento.
\item \codigoInline{portaAberta} é comportamento; \codigoInline{abrirPorta()} é estado.
\item Ambos são estados porque pertencem ao mesmo objeto.
\item Ambos são comportamentos porque podem ser acessados com ponto.
\item O campo é uma classe e o método é um objeto.
\end{alternativas}
\questaoDireta{6}{Queremos um método que receba um andar de destino inteiro e não retorne valor. Qual assinatura atende ao requisito?}{5 pontos}
\begin{alternativas}
\item \codigoInline{int irPara()}
\item \codigoInline{void irPara()}
\item \codigoInline{int irPara(int destino)}
\item \codigoInline{void irPara(int destino)}
\item \codigoInline{void irPara(float portaAberta)}
\end{alternativas}
\questaoDireta{7}{O código compila? Se não, qual linha impede a compilação?}{5 pontos}
\begin{painelEvolucao}{Trecho em Elevador.java}
\begin{lstlisting}[style=cenariocodigo]
int andarAtual;
void mostrar() {
    int anterior;
    System.out.println(andarAtual);
    System.out.println(anterior);
}
\end{lstlisting}
\end{painelEvolucao}
\begin{alternativas}
\item Compila e as linhas 4 e 5 imprimem \codigoInline{0}.
\item Não compila: a linha 5 tenta ler a variável local \codigoInline{anterior} sem inicializá-la.
\item Não compila: a linha 4 tenta ler o campo \codigoInline{andarAtual} sem inicializá-lo.
\item Compila, mas a execução falha na linha 5 ao acessar \codigoInline{anterior}.
\item Não compila: apenas declarar a variável local na linha 3 já é inválido.
\end{alternativas}
\end{minipage}\hfill
\begin{minipage}[t]{.485\textwidth}
\questaoDireta{8}{O elevador suporta no máximo 8 pessoas. Ao embarcar um grupo, a lotação não pode ultrapassar esse limite. Qual solução mantém essa regra em um único responsável?}{5 pontos}
\begin{alternativas}
\item Cada trecho de \codigoInline{Main} consulta a lotação e decide se soma o grupo.
\item Um auxiliar de \codigoInline{Main} calcula a lotação e altera o campo.
\item \codigoInline{Main} verifica o grupo; outro método externo verifica o limite.
\item Um setter recebe a lotação calculada pelo código externo.
\item \codigoInline{Elevador.embarcar(quantidade)} consulta sua lotação e preserva o limite.
\end{alternativas}
\questaoDireta{9}{Considere o método e a chamada abaixo. O que acontece?}{5 pontos}
\begin{painelEvolucao}{Trecho em Elevador.java e Main.java}
\begin{lstlisting}[style=cenariocodigo]
int calcularDiferenca(int destino) {
    return destino - andarAtual;
}
int diferenca = elevador.calcularDiferenca(3);
\end{lstlisting}
\end{painelEvolucao}
Considere que \codigoInline{andarAtual} vale \codigoInline{1} antes da chamada.
\begin{alternativas}
\item \codigoInline{diferenca} recebe \codigoInline{3} e \codigoInline{andarAtual} passa a ser \codigoInline{3}.
\item \codigoInline{diferenca} recebe \codigoInline{2} e \codigoInline{andarAtual} passa a ser \codigoInline{3}.
\item O código não compila porque um método com retorno \codigoInline{int} não pode receber parâmetro.
\item \codigoInline{diferenca} recebe \codigoInline{2} e \codigoInline{andarAtual} continua em \codigoInline{1}.
\item O método não retorna valor porque não foi declarado com \codigoInline{void}.
\end{alternativas}
\questaoDireta{10}{Dois elevadores são criados separadamente e começam no térreo com a porta fechada. Depois, apenas \codigoInline{a.subir()} é chamado. Qual afirmação está correta?}{5 pontos}
\begin{alternativas}
\item Antes da chamada, \codigoInline{a == b} é \codigoInline{true}; depois, ambos ficam no andar 1.
\item Antes da chamada, os estados diferem; depois, somente \codigoInline{b} fica no andar 1.
\item Antes da chamada, os estados são iguais e \codigoInline{a == b} é \codigoInline{false}; depois, somente \codigoInline{a} fica no andar 1.
\item Antes da chamada, \codigoInline{a == b} é \codigoInline{true}; depois, somente \codigoInline{a} fica no andar 1.
\item Antes da chamada, os estados são iguais; por isso, alterar \codigoInline{a} também altera \codigoInline{b}.
\end{alternativas}
\questaoDireta{11}{\codigoInline{getAndarAtual()} não possui modificador de acesso. Como \codigoInline{Main} e \codigoInline{Elevador} estão no mesmo pacote, qual afirmação está correta?}{5 pontos}
\begin{alternativas}
\item \codigoInline{Main} pode chamá-lo: o acesso de pacote permite uso no mesmo pacote.
\item \codigoInline{Main} não pode chamá-lo porque o método não é \codigoInline{public}.
\item Ele só pode ser chamado se \codigoInline{andarAtual} também for público.
\item Ele se torna privado automaticamente por não declarar \codigoInline{public}.
\item Ele altera o pacote de \codigoInline{Main} durante a execução.
\end{alternativas}
\questaoDireta{12}{O elevador só pode ficar entre os andares 0 e 3. Analise a alteração abaixo. Qual diagnóstico está correto?}{5 pontos}
\begin{painelEvolucao}{Trecho em Elevador.java}
\begin{lstlisting}[style=cenariocodigo]
private int andarAtual;
public void definirAndar(int novoAndar) {
    andarAtual = novoAndar;
}
\end{lstlisting}
\end{painelEvolucao}
\begin{alternativas}
\item O campo \codigoInline{private} garante sozinho que todo andar seja válido.
\item Receber um \codigoInline{int} impede automaticamente andares negativos.
\item Nem os métodos da própria classe podem alterar o campo.
\item O acesso direto fechou, mas o método ainda aceita qualquer andar.
\item Um método público não pode alterar um campo privado.
\end{alternativas}
\end{minipage}

\end{document}
